% =========================================================================== %
% paper-en.tex —— English top-venue authoring template (A-track simulation paper)
%
% BUILD
%   lualatex paper-en.tex && lualatex paper-en.tex
%   (LuaLaTeX, not XeLaTeX: microtype font expansion needs LuaTeX and is the
%    largest single contributor to the "professionally typeset" impression.)
%
% HOW TO USE THIS FILE
%   The section order below IS the argument order. Do not reorder it. A top-venue
%   paper is not "restate the problem, then analyse it, then model it"; it is
%   "here is a gap, here is a mechanism, here is a method that follows from the
%   mechanism, here is the evidence, here is where it breaks".
%   The moves for each section are specified in
%   `references/english-narrative.md`; `scripts/csf_prose.py` machine-checks the
%   subset that can be checked.
%
%   Write ENGLISH FIRST. Localisation to Chinese is a mechanical derivative of a
%   frozen English source (`scripts/csf_localize.py`), never a rewrite. Writing
%   Chinese first and translating is what produced the flat voice this template
%   exists to prevent.
%
% SWITCHING VENUE APPEARANCE
%   \usepackage[neurips,stix,final]{csfstyle-en}   (default: NeurIPS-like, 1 col)
%   \usepackage[icml,stix,final]{csfstyle-en}
%   \usepackage[nature,stix,final]{csfstyle-en}    (unnumbered sans headings)
%   \usepackage[aamas,charter,final]{csfstyle-en}  + [twocolumn] on the class
%   Add `draft` to show line numbers and make \csfclaim / \csfevi visible.
% =========================================================================== %

\documentclass[10pt]{article}
\usepackage[neurips,stix,final]{csfstyle-en}

% --- Bibliography. `plainnat` + natbib gives author-year, which is what every
% --- venue in scope uses. Swap for the venue .bst from _vendor/venue-styles/.
\usepackage[round,authoryear]{natbib}
\bibliographystyle{plainnat}

% --- Graphics. Vector only. A raster figure in a simulation paper is a defect:
% --- it cannot be re-scaled, and reviewers zoom.
\usepackage{graphicx}
\graphicspath{{figures/}{figs/}}

\begin{document}

% --------------------------------------------------------------------------- %
% Title block. The title states the FINDING or the MECHANISM, not the topic.
%   weak   : "Research on Optimization of Multi-Agent Emergency Evacuation"
%   strong : "Capacity-Proportional Allocation Beats Congestion-Aware Routing in
%             Multi-Agent Evacuation: A Mechanism-Level Account"
% The weak form announces a topic and promises nothing; it is the single most
% common title defect in contest drafts.
% --------------------------------------------------------------------------- %
\csftitle{<Finding-first title: state the mechanism or the result, not the topic>}

\csfauthors{<Author One$^{1}$, Author Two$^{1,2}$>}
\csfaffil{$^{1}$<Institution> \quad $^{2}$<Institution>}
\csfemail{\{author.one, author.two\}@example.edu}

\begin{csfabstract}
% Five moves, 4-6 sentences total, no citations, no equations, no "\ref".
%   (1) Context in one clause, with the stakes and, if possible, a number.
%   (2) The gap, stated as a mechanism that is missing or misunderstood.
%   (3) What we do: the model and the one design decision that matters.
%   (4) The quantitative result, with the comparison and the interval.
%   (5) What it implies, in one clause, without hedging.
% Each sentence that makes a claim must be supported somewhere in the body; an
% abstract claim with no body evidence is an automatic reviewer objection.
<TODO move 1: context and stakes of the simulated system, in one clause.>
<TODO move 2: the specific gap, phrased as a mechanism, not as "few studies".>
<TODO move 3: our model and the single design decision that drives the result.>
<TODO move 4: headline number, with baseline and uncertainty.>
<TODO move 5: the implication, stated plainly.>

\csfkeywords{multi-agent simulation; <domain>; <method family>; <mechanism>; reproducibility}
\end{csfabstract}

% --------------------------------------------------------------------------- %
\section{Introduction}
% Five paragraphs, in this order. The funnel is not decoration: it is the
% sequence a reviewer needs in order to accept the contribution.
%
%   P1 Domain and stakes, with at least one real number and a citation.
%   P2 The specific problem, narrowed to what this paper solves.
%   P3 Why the obvious approaches fail, AT THE LEVEL OF MECHANISM. This is the
%      paragraph that makes the paper scientific rather than engineering.
%   P4 Our approach in two sentences plus the intuition for why it should work.
%   P5 Contributions, as a numbered list, each ending in a testable claim.
%
% The contribution list is the contract with the reader. Every C-item must be
% discharged by evidence later, and csf_prose.py checks that each one names the
% figure, table or section that discharges it.
% --------------------------------------------------------------------------- %

<TODO P1: domain and stakes. One number. One citation, e.g. \citep{helbing1995social}.>

<TODO P2: the specific problem, narrowed. Define the decision that has to be made.>

<TODO P3: why the obvious approach fails, by mechanism. Not "existing methods have
limitations" but "routing on instantaneous queue length is self-defeating because
the act of choosing a station changes the quantity being chosen on".>

<TODO P4: our approach, and the intuition that motivates it.>

\begin{csfcontrib}
\begin{csfcontribs}
  \item <A formal model of ... (discharged by \cref{sec:model})>
  \item <A mechanism-level explanation of why ... (discharged by \cref{fig:mechanism})>
  \item <An allocation rule that ... (discharged by \cref{tab:main}, \cref{fig:main})>
  \item <A cross-implementation parity check that ... (discharged by \cref{sec:repro})>
\end{csfcontribs}
\end{csfcontrib}

% --------------------------------------------------------------------------- %
\section{Related Work}
% Structure by CONTRAST, never as an annotated bibliography. Each paragraph takes
% one family of prior work and ends with the specific difference that matters:
%   "X et al. optimise Y under Z. Their formulation assumes <assumption>, which
%    fails in <our setting> because <mechanism>. We therefore <difference>."
% A paragraph that only summarises is dead weight; delete it.
% csf_prose.py fails the paper when NO paragraph contains a contrast marker.
% --------------------------------------------------------------------------- %

<TODO: 3-5 contrast paragraphs, each ending in a concrete difference.>
% Cite prior work with \citet{...} where the authors are the subject of the
% sentence and \citep{...} where the work is parenthetical. Example:
%   "Social-force models reproduce local interactions \citep{helbing1995social}
%    but assume exits are equivalent, whereas our formulation makes the allocation
%    decision endogenous."
% Keep the two example lines above COMMENTED: leaving a literal \citet{...} in
% live text sends the key "..." to BibTeX, which then warns that it found no
% database entry for it. That warning is noise that hides real ones.

% --------------------------------------------------------------------------- %
\section{Problem Formulation}
\label{sec:problem}
% The formalisation must be a definition of the DECISION, not a notation dump.
% Order: system description -> state/action/observation -> objective -> what makes
% it hard. State the objective as an optimisation with explicit constraints; an
% objective without constraints is not a formulation.
% The notation table is mandatory and must be referenced before it appears.
% --------------------------------------------------------------------------- %

<TODO: system definition, agent set, state, action, observation, transition.>

<TODO: objective and constraints. Give the lower bound / best case if one exists —
a bound turns every later number into a ratio and makes the results interpretable.>

% NOTE ON FLOAT WIDTH — this matters as soon as you switch to a two-column venue.
% A data table with four columns does NOT fit an ICML/AAAI column (~3.2in), and
% LaTeX reports it as an Overfull \hbox rather than failing: the table silently
% runs into the gutter. Verified by compiling this exact template with
% [twocolumn,icml,...]: both tables below produced Overfull \hbox (16.5pt and
% 55.4pt too wide).
% `table*` spans both columns, and in a single-column document it behaves exactly
% like `table`, so the starred form is the portable choice for wide data tables.
% Use the unstarred form only for narrow tables that genuinely fit one column.
\begin{table*}[t]
  \centering
  \caption{Notation. Every symbol used in \cref{sec:model} must appear here, and
    no symbol that appears only once in the paper should.}
  \label{tab:notation}
  \begin{tabular}{@{}lll@{}}
    \toprule
    Symbol & Meaning & Domain \\
    \midrule
    $N$ & number of agents & $\N$ \\
    $S$ & set of service stations & --- \\
    $c_j$ & capacity (service rate) of station $j$ & $\R_{>0}$ \\
    $T$ & makespan (last completion time) & $\R_{\ge 0}$ \\
    $T^{\mathrm{lb}}$ & capacity lower bound on $T$ & $\R_{\ge 0}$ \\
    \bottomrule
  \end{tabular}
\end{table*}

\begin{csfasm}[Assumptions]
% Each assumption needs: the statement, its real-world justification, and what
% changes if it is relaxed. An assumption list without the third item is a list of
% excuses; the third item is what a reviewer is actually looking for.
\begin{enumerate}
  \item <Statement.> \emph{Justification:} <why it holds in this domain.>
    \emph{If relaxed:} <which result changes and in which direction.>
\end{enumerate}
\end{csfasm}

% --------------------------------------------------------------------------- %
\section{Model and Method}
\label{sec:model}
% For every component, give the design RATIONALE, not just the design. The pattern
% to use is: "Because <mechanism>, the allocation must satisfy <property>, which
% rules out <obvious choice> and leaves <our choice>." A component presented
% without its why is indistinguishable from a component copied from a textbook.
% Use run-in \paragraph{...} headings for these rationale blocks; they read far
% better than another numbered level.
% --------------------------------------------------------------------------- %

\subsection{Model}
<TODO: transitions, decision epochs, information available to each agent, and the
crucial question of WHEN an agent may re-decide.>

\subsection{Allocation rule}
<TODO: state the rule, then prove or argue the property it has.>

\begin{csfproposition}
  <A property of the rule, stated so it can be falsified.>
\end{csfproposition}
\begin{proof}
  <Proof, or, for an empirical claim, an explicit pointer to the experiment that
  tests it. Do not dress up a simulation observation as a theorem.>
\end{proof}

\paragraph{Design rationale.} <Why this rule and not the two obvious alternatives.>

\subsection{Algorithm}
<TODO: pseudocode only if the procedure cannot be stated in prose; a 40-line
listing that restates the equations is padding.>

% --------------------------------------------------------------------------- %
\section{Experiments}
\label{sec:exp}
% Order: setup -> baselines -> main result -> mechanism -> ablation -> sensitivity
%        -> failure cases.
%
% The baseline TAXONOMY matters more than the baseline count. A reviewer wants to
% see at least three distinct classes represented: a rule/heuristic, a
% single-agent learned method, and a multi-agent learned method. Three variants of
% one family is not a comparison, and \code{csf\_readiness.py} fails the paper for it.
% --------------------------------------------------------------------------- %

\subsection{Setup}
<TODO: instance generation, size, seeds, hardware, wall-clock. State the number of
independent runs and the seed policy — "we report the mean over 10 seeds" without
a seed policy is not reproducible.>

\subsection{Baselines}
<TODO: one sentence per baseline, naming its class and its information access. A
baseline given less information than the proposed method is not a baseline.>

\subsection{Main result}
% The caption must state the comparison rule: what is bold, what the error bars
% are, how many runs. Never bold the best value without stating the rule, and
% never claim a win that is inside the noise.
% `table*` for the same reason as the notation table above: four columns of data
% overflow a two-column measure silently (Overfull \hbox, not an error).
\begin{table*}[t]
  \centering
  \caption{Main comparison on <instance family>. Mean $\pm$ s.d.\ over $n=<k>$
    seeds. Bold marks the best mean within one standard deviation;
    $T^{\mathrm{lb}}$ is the capacity lower bound from \cref{sec:problem}.}
  \label{tab:main}
  \begin{tabular}{@{}lS[table-format=3.2]S[table-format=1.3]l@{}}
    \toprule
    Method & {$T$ (min)} & {$T/T^{\mathrm{lb}}$} & Class \\
    \midrule
    Nearest station            & 483.94 & 3.387 & rule \\
    Random assignment          & 238.56 & 1.670 & rule \\
    Capacity-proportional (\ours{ours}) & \best{157.01} & \best{1.099} & rule + bound \\
    \midrule
    Capacity lower bound       & 142.86 & 1.000 & --- \\
    \bottomrule
  \end{tabular}
\end{table*}

<TODO: one paragraph that reads the table. State the effect size, the direction,
and the mechanism that explains it. "Our method performs better" is not a reading
of a table; "the gap closes from 3.4x to 1.1x of the lower bound, which means the
remaining inefficiency is no longer allocation but travel" is.>

\begin{csftake}
  <One sentence the reader should be able to repeat from memory after closing the
  paper. If you cannot write this sentence, the experiment does not have a result
  yet.>
\end{csftake}

\subsection{Mechanism}
% This subsection is what separates a simulation paper from a benchmark report:
% show WHY the effect exists, not only that it does. Plot the intermediate
% quantity, not just the outcome.
<TODO: the intermediate quantity that carries the effect.>

\subsection{Ablation}
<TODO: one row per design decision, each removed individually. An ablation that
changes two things at once explains nothing.>

\subsection{Sensitivity and robustness}
<TODO: vary the parameters the conclusion depends on, and report the range over
which the conclusion survives. "Robust" must come with the range in which it is
robust.>

\subsection{When it fails}
<TODO: the regime where the method loses, with the numbers that show it. A paper
with no failure case is either trivial or not yet understood.>

% --------------------------------------------------------------------------- %
\section{Discussion}
<TODO: what the mechanism implies beyond this instance, and what would falsify the
mechanism. Keep it short; speculation padded to a page is worse than one honest
paragraph.>

% --------------------------------------------------------------------------- %
\section{Limitations}
% Mandatory and non-negotiable. Concrete limits, not "our method has limitations".
% Each limitation names the assumption it comes from and the direction of the
% resulting bias.
<TODO: three concrete limits with the direction of the bias each induces.>

\section*{Reproducibility Statement}
\csfstatement{Reproducibility.}
<TODO: where the code, instances and seeds live; which command reproduces which
figure; the exact software versions. A statement that says "code available on
request" is worth nothing and \code{csf\_readiness.py} reports it.>

\section*{Compute Statement}
\csfstatement{Compute.}
<TODO: hardware, wall-clock, number of runs. Required by ICML and increasingly by
everyone else.>

% --------------------------------------------------------------------------- %
% Figures. Place them near the text that needs them, and never refer to a figure
% without saying what the reader should conclude from it.
%
% DRAFTING BEFORE THE ART EXISTS: \csffigplaceholder occupies the EXACT final size
% and carries the caption, so line breaking, page count and caption length are all
% final before any drawing is done. Replace it with \includegraphics when the
% vector art is ready.
% --------------------------------------------------------------------------- %
\begin{figure}[t]
  \centering
  \csffigplaceholder{4.2cm}
  \caption{<What the reader must conclude, not what the figure contains.>
    Panels: (a)~<role>; (b)~<role>. Labels and vector art are produced separately;
    see \texttt{figures/contracts/} for the label ledger.}
  \label{fig:main}
\end{figure}

\bibliography{refs}

\end{document}
